\documentclass[10pt]{article}
\usepackage[a4paper,margin=2.6cm]{geometry}
\usepackage{array}
\usepackage{booktabs}
\usepackage{amsmath}
\usepackage{amssymb}
\usepackage[hidelinks]{hyperref}
\usepackage{xcolor}
\setlength{\parskip}{2pt}
\newcommand{\code}[1]{\texttt{#1}}

\title{Family 7.1: ocean colour joins the global tensor\\
\large a fourth channel group, inherited rather than rebuilt}
\author{blauewelt/earth --- design note}
\date{13 September 2026}

\begin{document}
\maketitle

\begin{abstract}
Family 7.1 extends the global input tensor of family 7 --- the Earth's
surface state on a $0.25^\circ$ grid from pole to pole in five-day bins,
1982--2024, in 54 channels --- with a fourth channel group carrying the
observed surface chlorophyll-a concentration from the ESA Ocean Colour
Climate Change Initiative (4\,km, September 1997 to December 2024). The
group holds two channels: the block mean of $\log_{10}$ chlorophyll over the
36 source cells of each $0.25^\circ$ point, and the fraction of cell-days in
the pentad that carried a retrieval. Its time axis starts where the record
starts, so the group is 1,997 bins rather than 3,142. The three inherited
groups are not recomputed: the build hard-links the published family-7
files, refuses any write through the new names, and its manifest states that
the inherited files hash identically to family 7's. We give the source, the
averaging rule, the missing-value semantics, the layout, the assertions, and
the reason the channel is a target before it is an input.
\end{abstract}

\section{Purpose}

The programme's forward model reads a per-cell state, predicts it forward,
and reads out quantities of interest from the predicted state. Surface colour
--- the chlorophyll-a concentration retrieved from water-leaving radiance ---
is the only global, routine observation of the ocean biosphere, and it is one
of the read-outs the programme has committed to (E-072). A read-out needs a
target on the tensor's own grid and time axis; family 7 had none. Family 7.1
adds one.

Two decisions frame the addition. First, colour enters as a \emph{target}
and probe before it enters as an input: the question it is built to answer is
whether the physical state the model already carries --- temperature, mixed
layer, currents, wind, light through the calendar --- predicts the biosphere a
season ahead. Second, the source is averaged onto the $0.25^\circ$ grid
rather than kept at its native 4\,km. That is a deliberate and temporary
simplification: it lets the channel drop into the three-group layout with no
new machinery, at the price of discarding a factor of 36 in pixels. The
companion note on family 10 describes the layout in which fine sources are
kept at their own granularity; family 7.1 is a tier-G group of that design,
unchanged.

\section{Source}

ESA Ocean Colour CCI version 6.0, geographic projection, daily, chlorophyll-a
only \cite{occci}: the merged SeaWiFS, MERIS, MODIS-Aqua, VIIRS and OLCI
record, bias-corrected across sensors, on a $1/24^\circ$ grid of
$4320 \times 8640$ cells with centres at odd multiples of $1/48^\circ$,
4~September~1997 to 31~December~2024. The variable is \code{chlor\_a} in
mg\,m$^{-3}$, float32 with a fill value. The retrieval exists only where the
sensor saw the water: not under cloud, sun glint or sea ice, not in polar
night, not over land. The builder does not assume file names; it lists each
year's directory on the archive's own catalogue and preflights one real file
before spending anything. The daily uncertainty fields shipped with the
product are not stored --- an uncertainty of a merged algorithm is not an
observation --- and the gap-filled level-4 products are not used, because the
target should be what was seen.

\section{Channels}

\begin{table}[h]
\centering\setlength{\tabcolsep}{4pt}\footnotesize
\begin{tabular}{@{}rl>{\raggedright\arraybackslash}p{2.2cm}>{\raggedright\arraybackslash}p{9.2cm}@{}}
\toprule
\# & channel & unit & rule \\
\midrule
0 & \code{log\_chl} & $\log_{10}$\,mg\,m$^{-3}$ & per day, the mean of $\log_{10}$ chlorophyll over the finite cells of the $6\times6$ block centred on the point; per pentad, the mean of those daily means over the days with at least one finite cell; absent if no day had one \\
1 & \code{chl\_cov} & fraction in $(0, 1]$ & finite cell-days in the block during the pentad, divided by (cells in the block $\times$ days in the bin); absent exactly where \code{log\_chl} is \\
\bottomrule
\end{tabular}
\caption{Group \code{oc025}: two ocean-colour channels at $0.25^\circ$, shape $[1997, 721, 1440, 2]$, z-scored per channel like every other group, the standardisation constants stored as \code{norm\_oc025}.}
\end{table}

\paragraph{The block.} A point of the $0.25^\circ$ grid sits at a multiple of
$0.25^\circ = 6/24^\circ$, which is a cell boundary of the $1/24^\circ$ grid;
the six cells whose centres lie within $\pm 0.125^\circ$ of the point on each
axis are therefore exactly three on either side. The block wraps in
longitude; at the two pole rows it is truncated to the cells that exist. The
implementation is a roll by three cells and a reshape, so the pole and the
dateline are the only special cases and both are exercised by tests.

\paragraph{Log space.} Chlorophyll is log-normal over four orders of
magnitude. The arithmetic mean of a block containing one bloom cell at
30\,mg\,m$^{-3}$ and thirty-five at 0.05 describes neither; the mean of
$\log_{10}$ does. The programme's read-out for colour is scored in log space
(E-072), so the stored quantity and the scored quantity agree.

\paragraph{Coverage.} A pentad mean resting on one clear pixel on one day and
one resting on thirty-six pixels on five days are different amounts of
evidence, and the model should be told which it is looking at.
\code{chl\_cov} carries that. It is stored un-logged and z-scored like every
other channel; its distribution is a property of the observing system (clear
skies at low latitude, cloud and polar night at high) and is reported from
the built file rather than assumed here.

\section{Missing values}

\code{NaN} keeps its family-7 meaning --- not observed here at this time ---
and for colour it is the rule rather than the exception: at high latitudes in
winter no day is observed for months. Before bin 1145, the bin that contains
4~September~1997, there is no observation at all, and the group's time axis
starts there. \code{chl\_cov} is absent exactly where \code{log\_chl} is; a
block that saw nothing has no coverage to report, and storing zero there would
collapse the distinction between ``unobserved'' and ``observed to be empty''
that the missing-token design relies on. Nothing is filled: no climatological
fill, no interpolation across days, no level-4 product.

\section{Layout and inheritance}

The group is one file, \code{family7\_global025\_pentad\_l1\_X\_oc025.npy}:
float16, C order, a 128-byte NumPy header, bin-major so that one pentad is a
contiguous 4.15\,MB range read. Its first axis is offset --- row $r$ holds
bin $1145 + r$ --- and the offset is stored in the metadata file as
\code{oc\_bin\_first} beside an explicit bin index. The offset saves 4.7\,GB
of empty slabs (8.3\,GB instead of 12.9); it costs a per-group first bin in
every consumer, which the samplers and the front end already supported for
the monthly Argo group.

Family 7.1 is a new recipe, \code{f7l1}, with its own stem and its own
folder on the public dataset repository, because the family-7 files, their
hashes and their documentation are cited elsewhere and a published dataset
should not change under its own name. The three inherited groups are
\emph{hard-linked} from the finished family-7 build directory --- one inode,
no extra bytes --- and never recomputed. Since a hard link shares its inode,
a write through the new name would corrupt the published bytes, and the
process runs as root, which ignores file modes; so the refusal is structural
rather than a permission bit: every writable open of a group file and the
stage machinery's discard path check whether the group was seeded and exit
if so. The publish step fetches family 7's manifest and asserts that each
inherited file hashes identically before it writes \code{same\_as\_f7l0:
true} beside that file's entry. The metadata file is regenerated and every
key it shares with family 7's is required to be equal.

\section{Assertions}

Before the tensor is trusted: the inherited files are byte-identical to
family 7's manifest; the group's shape, dtype and header length are as
stated; the first bin is computed from the epoch, not typed; every finite
coverage lies in $(0, 1]$ and \code{log\_chl} is finite exactly where
coverage is; the block mean reproduces exact values on synthetic input (a
block of $10^{0.5}$ on all five days reads $0.5$ with coverage $1$; one
finite cell on one day reads its own $\log_{10}$ with coverage $1/180$); the
pole rows and the dateline are exercised; the stage resumes across a
simulated interruption; and on the bin opening 3~January~2015 --- the
comparison bin of the family-7 validation --- the global mean of
\code{log\_chl} and the fraction of observed ocean cells are written to the
build record for comparison, by eye, against the chlorophyll layers already
on the globe.

\section{Holdout}

Unchanged: 2009, 2017 and 2023 are development holdout years; the terminal
split trains through 2020 and tests on 2021--2024. Colour's record begins in
September 1997, so a training-years-only seasonal climatology for the channel
has twenty-one years to fit from.

\section{Limitations}

The averaging discards 35 of every 36 source pixels; the coverage channel
records the loss but does not undo it. The channel is clear-sky and daylight
only, which halves its live bins at high latitude and removes the polar
winter entirely. Reflectance bands, which the programme's rule names as the
input from which colour is the target, are not carried; they belong to the
first design that keeps fine sources at their own granularity. And the record
is 27 years against the tensor's 43.

\begin{thebibliography}{9}\small
\bibitem{occci} Sathyendranath, S.\ et al.\ (2019). An Ocean-Colour Time
Series for Use in Climate Studies: The Experience of the Ocean-Colour Climate
Change Initiative (OC-CCI). \emph{Sensors} 19(19), 4285. Version 6.0 data:
ESA Climate Office / Plymouth Marine Laboratory.
\bibitem{f7} blauewelt/earth (2026). Family 7: a global, three-resolution
input tensor for a forward model of the Earth's surface state. Design note,
7~September~2026.
\end{thebibliography}

\end{document}
